% !PRINTED 705
\weberpage{705}
Diese Funktion muß also in $\mathfrak P_2, \mathfrak P_3, \ldots$ verschwinden. In dem Punkte $\mathfrak P_2$ verschwindet aber $z\,x_{1,1}\varrho_1, z\,x_{1,3}\varrho_3, \ldots, z\,x_{1,n}\varrho_n$, während $\varrho_2$ nicht verschwindet. Folglich muß $z\,x_{1,2}$ in $\mathfrak P_2$, d.\,h. für $z=\infty$, verschwinden und $z^2 x_{1,2}$ endlich bleiben. Wir haben also:

\gesperrt{Ist $r$ nicht gleich $s$, so ist $z^2 x_{r,s}$ für $z=\infty$ endlich.}

Wir definieren nun die rationale Funktion $u^{(i)}$ durch:
\weq{11}{x_{i,i} = a^{(i)}_{m-2} z^{m-2} + a^{(i)}_{m-3} z^{m-3} + \cdots + a^{(i)}_{-1} z^{-1} + u^{(i)} z^{-2},}
worin die Konstante $a^{(i)}_\nu$ dieselbe sein soll wie in (5), und wir erhalten:
\weq{12}{\frac{dJ}{dz} - x_{i,i} = z^{-2}\bigl[\eta^{(i)} - u^{(i)}\bigr],}
und nach (8):
\[
\bigl[\eta^{(i)} - u^{(i)}\bigr]\varrho_i = z^2 x_{i,1}\varrho_1 + \cdots + z^2 x_{i,i-1}\varrho_{i-1} + z^2 x_{i,i+1}\varrho_{i+1} + \cdots + z^2 x_{i,n}\varrho_n.
\]

Im Punkte $\mathfrak P_i$ ist $\eta^{(i)}$ endlich und $\varrho_i$ von Null verschieden, während auf der rechten Seite alles in $\mathfrak P_i$ verschwindet. Folglich muß auch $u^{(i)}$ in $\mathfrak P_i$, und, weil es rational ist, für $z=\infty$ endlich sein, und aus (9) und (11) folgt:
\weq{13}{S\left(\frac{dJ}{dz}\right) = \overset{i}{\Sigma} a^{(i)}_{m-2} z^{m-2} + \cdots + \overset{i}{\Sigma} a^{(i)}_{-1} z^{-1} + \overset{i}{\Sigma} u^{(i)} z^{-2}.}

Nach unserer Annahme über $z$ enthält das Untereck $\mathfrak B$ von $dJ$ keinen Punkt, in dem $z$ einen endlichen Wert hat.

Ist $\mathfrak P$ ein Punkt, in dem $z$ den endlichen Wert $z_0$ hat, so ist
\weq{14}{(z - z_0) = \frac{\mathfrak N}{\mathfrak U}, \quad \frac{dJ}{dz} = \frac{\mathfrak A \mathfrak U^2}{\mathfrak Z\mathfrak B}, \quad (z-z_0)\frac{dJ}{dz} = \frac{\mathfrak U \mathfrak N \mathfrak A}{\mathfrak Z\mathfrak B},}
und in $\mathfrak N$ ist der Punkt $\mathfrak P$ einmal öfter enthalten als in $\mathfrak Z$, folglich ist:
\weq{15}{\left[(z-z_0)\frac{dJ}{dz}\right]_0 = 0 \quad \text{(in $\mathfrak P$).}}
Das ist die Forderung a) in \S\,198.

Ist daher $\lambda_1, \lambda_2, \ldots, \lambda_n$ die Normalbasis in $z$, und $\mu_1, \mu_2, \ldots, \mu_n$ die komplementäre Basis, so ist:
\weq{16}{\frac{dJ}{dz} = y_1 \mu_1 + y_2 \mu_2 + \cdots + y_n \mu_n,}
worin die $y_i$ ganze rationale Funktionen von $z$ sind, woraus folgt, daß
\weq{17}{S\left(\frac{dJ}{dz}\right) = y_1 \quad [\S\,198, (18)]}
\gesperrt{eine ganze rationale Funktion von $z$ ist.}
